% SPDX-FileCopyrightText: Copyright (c) 2024-2026 Yegor Bugayenko
% SPDX-License-Identifier: MIT

\section{File System}\label{sec:fs}

Here we define manipulations with files and directories.
All objects presented in this Section belong to \ff{fs} package.

The object \deff{path} is an abstraction of file path and is a decorator of \deff{string}.
%
The attribute \deff{separator} is a system-dependent default name-separator character.
%
The attribute \deff{joined} is a utility object that joins the given \deff{tuple}
  of paths with the current OS separator and normalizes the result path.
%
The attribute \deff{as-file} is a \deff{file} created from the path itself.
%
The attribute \deff{as-dir} is a \deff{dir} created from the path itself.
%
The attribute \deff{is-absolute} is \deff{true} if the path is absolute and \deff{false} otherwise.
%
The attribute \deff{normalized} is a new \deff{path} with normalized \(uri\).
Normalization includes converting multiple slashes into a single slash
  and resolving ``.'' (current directory) and ``..'' (parent directory) segments.
%
The attribute \deff{resolved} is a new \deff{path} with a suffix appended to the current one.
%
The attribute \deff{basename} is the name of the file in the path.
%
The attribute \deff{extname} is the extension in the path.
%
The attribute \deff{dirname} is the name of the directory in the path.

The object \deff{file} is an abstraction of a file and is a decorator of \deff{path},
  which is the path of the file.
%
The attribute \deff{as-path} is a \deff{path} of the file.
%
The attribute \adeff{is-directory} is \deff{true} if the file is a directory.
%
The attribute \adeff{exists} is \deff{true} if the file exists.
%
The attribute \deff{touched} makes sure the file exists.
%
The attribute \deff{deleted} removes the file.
%
The attribute \adeff{size} is the size of the file in bytes.
%
The attribute \deff{moved} renames or moves the file to a new place.
%
The attribute \deff{open} opens the file.
The \deff{open} object accepts two attributes: file open mode which is a \deff{string}
  and file scope.
The file scope attribute must be an abstract object with one void attribute
  which is \deff{file-stream}.
It has the attributes \adeff{read} and \adeff{write} which allow working with the file
  as ``input'' and ``output''.
The result of dataization of \ff{file.open} is the result of dataization
  of its second argument ``scope''.
The file descriptor is closed automatically when the scope is dataized.

The object \deff{dir} is an abstraction of a directory.
%
The attribute \deff{made} creates the directory with all its parent directories.
%
The attribute \deff{deleted} deletes the directory recursively with all its subdirectories.
%
The attribute \adeff{walk} finds files in the directory using glob pattern matching.
%
The attribute \deff{tmpfile} is a temporary file in the directory.

The object \deff{tmpdir} is a decorator of \deff{dir} that abstracts the system
  directory for temporary files.

In the next example, a temporary file is created and filled up with \ff{"Hello, world"}.
Then the data is read from the file and written to \deff{malloc}.

\begin{ffcode}
malloc.of
  12
  seq * > [m]
    tmpdir
    .tmpfile
    .open
      "w"
      f.write "Hello, world" > [f]
    .open
      "r"
      m.put > [f] >>
        f.read 12
    m
\end{ffcode}
